% Standalone problem document, generated from the adjacent README.md.
% Compile with XeLaTeX. Mathematical content is maintained in Markdown.
\documentclass[11pt,a4paper]{article}
\usepackage[fontsize=10.5pt]{scrextend}
\usepackage[margin=22mm,headheight=15pt,headsep=9mm,footskip=11mm]{geometry}
\usepackage{fontspec}
\setmainfont{texgyrepagella}[Extension=.otf,UprightFont=*-regular,BoldFont=*-bold,ItalicFont=*-italic,BoldItalicFont=*-bolditalic]
\setsansfont{texgyreheros}[Extension=.otf,UprightFont=*-regular,BoldFont=*-bold,ItalicFont=*-italic,BoldItalicFont=*-bolditalic]
\setmonofont{texgyrecursor}[Extension=.otf,UprightFont=*-regular,BoldFont=*-bold,ItalicFont=*-italic,BoldItalicFont=*-bolditalic,Scale=MatchLowercase]
\usepackage{amsmath,amssymb}
\usepackage{unicode-math}
\setmathfont{texgyrepagella-math.otf}
\usepackage{xcolor,microtype,enumitem,needspace,fancyhdr}
\usepackage{bookmark}
\definecolor{ink}{HTML}{17344B}
\definecolor{accent}{HTML}{197D80}
\definecolor{muted}{HTML}{596773}
\hypersetup{colorlinks=true,linkcolor=accent,urlcolor=accent,pdftitle={MI-01 - Audenaert's
two-row Schatten norm-compression
conjecture},pdfauthor={Open Problems in Numerical Linear Algebra}}
\urlstyle{same}
\setlength{\parindent}{0pt}
\setlength{\parskip}{5pt plus 1pt minus 1pt}
\linespread{1.04}
\setlength{\emergencystretch}{3em}
\widowpenalty=10000
\clubpenalty=10000
\displaywidowpenalty=10000
\allowdisplaybreaks[1]
\setlist{nosep,leftmargin=1.5em}
\providecommand{\tightlist}{\setlength{\itemsep}{0pt}\setlength{\parskip}{0pt}}
\providecommand{\pandocbounded}[1]{#1}
\pagestyle{fancy}
\fancyhf{}
\fancyhead[L]{\sffamily\footnotesize\color{muted}OPEN PROBLEMS IN NUMERICAL LINEAR ALGEBRA}
\fancyhead[R]{\sffamily\bfseries\footnotesize\color{accent}MI-01}
\fancyfoot[L]{\parbox[b]{0.9\textwidth}{\sffamily\fontsize{8}{10}\selectfont\color{muted}Editorial ratings; a bounded search cannot certify that a problem remains unsolved.\\Literature check: 2026-09-08}}
\fancyfoot[R]{\sffamily\footnotesize\color{muted}\thepage}
\renewcommand{\headrulewidth}{0.3pt}
\renewcommand{\headrule}{\hbox to\headwidth{\color{accent}\leaders\hrule height \headrulewidth\hfill}}
\makeatletter
\renewcommand{\section}{\Needspace{5\baselineskip}\@startsection{section}{1}{0pt}{12pt plus 2pt minus 2pt}{4pt}{\sffamily\bfseries\large\color{ink}}}
\renewcommand{\subsection}{\Needspace{4\baselineskip}\@startsection{subsection}{2}{0pt}{9pt plus 1pt minus 1pt}{3pt}{\sffamily\bfseries\normalsize\color{ink}}}
\makeatother
\setcounter{secnumdepth}{0}
\begin{document}
{\sffamily\small\color{muted}Matrix inequalities and norms\par}
\vspace{3pt}
{\raggedright\hyphenpenalty=10000\exhyphenpenalty=10000\sffamily\bfseries\fontsize{21}{24}\selectfont\color{ink}Audenaert's
two-row Schatten norm-compression conjecture\par}
\vspace{7pt}
{\sffamily\small\textbf{Difficulty:} extreme\quad\textbf{Importance:} broadly
interesting\par}
\vspace{4pt}
{\color{accent}\hrule height 0.8pt}
\vspace{6pt}
\subsection{Problem statement}\label{problem-statement}

Let \(N,r_1,r_2,c_1,\ldots,c_N\) be positive integers and let
\(A_j\in\mathbb C^{r_1\times c_j}\),
\(B_j\in\mathbb C^{r_2\times c_j}\). For \(1\le p<\infty\), let
\(\|X\|_p=(\sum_i\sigma_i(X)^p)^{1/p}\), where \(\sigma_i(X)\) are the
singular values. Define

\[T=\begin{bmatrix}A_1&\cdots&A_N\\B_1&\cdots&B_N\end{bmatrix},\qquad
C_p(T)=\begin{bmatrix}\|A_1\|_p&\cdots&\|A_N\|_p\\\|B_1\|_p&\cdots&\|B_N\|_p\end{bmatrix}.\]

Do the following inequalities hold for every such partition?

\[\|T\|_p\ge\|C_p(T)\|_p\quad(1\le p\le2),\qquad
\|T\|_p\le\|C_p(T)\|_p\quad(2\le p<\infty).\]

\subsection{Why it matters}\label{why-it-matters}

The question asks whether the norm of a large block matrix can be
controlled sharply using only the norms of its blocks. Such compression
estimates are useful when analyzing matrix approximation errors without
retaining the full matrices.

\subsection{References}\label{references}

\begin{enumerate}
\def\labelenumi{\arabic{enumi}.}
\tightlist
\item
  K. M. R. Audenaert, \emph{On a norm compression inequality for
  \(2\times N\) partitioned block matrices}, arXiv:math/0702186v2 (26
  March 2007), Conjecture 1, equations (1)--(2), §§2 and 4. Published in
  Linear Algebra and its Applications 428(4) (2008).
  \href{https://arxiv.org/html/math/0702186}{Primary text}.
\item
  T. Zhang, \emph{Clarkson--McCarthy type inequalities, part I:
  \(\ell_p\)--\(\ell_p\) and \(\ell_q\)--\(\ell_p\) Schatten
  \(p\)-estimates}.\\
  arXiv:2410.21961v4 (28 March 2026), Conjecture 1.14 and Theorem 1.15.
  \href{https://arxiv.org/html/2410.21961}{Primary text}.
\end{enumerate}

\subsection{Status check --- 2026-09-08}\label{status-check-2026-09-08}

The 2026 revision still states this conjecture and supplies a weaker
bound. Audenaert proves several restricted cases, including \(p\ge4\);
the three-row generalization has counterexamples. Hanner's inequality, a
special case, has subsequently been proved, which does not settle
arbitrary two-row compression. Searches included
\textit{Audenaert norm compression conjecture 2026},
\textit{two row Schatten compression proof}, and
\textit{norm compression counterexample 2025 2026}; no general
resolution was located. Latest arXiv versions were checked for both
references.
\end{document}
